\section{Results}
\label{sec:results}
All 19 frozen cells completed 100 regenerated/refitted trials each, with 2,000 tests per trial: 44,600 method/group records, including 2,000 explicit unsmoothed-ILR inapplicability records. No pilot enters confirmation. Brackets below are pointwise Student-$t$ 95\% intervals over independent trials, with pairing for contrasts; they are not simultaneous theorem tests or guarantees for each calibration sample. Appendix~\ref{app:full-results} retains every declared method/group entry, domain outcome and denominator; Appendix~\ref{app:measured-regions} shows all 32 prespecified measured cases.
\input{figures/core_main}

\paragraph{Q1: coordinate calibration and joint reliability.}
Naive coordinate calibration gives interior joint coverage $.7806$ [$ .7766,.7847$] at $K=3$ and $.3215$ [$ .3176,.3254$] at $K=10$ (Table~\ref{tab:core_coordinate_joint}). In the latter cell the ten coordinate means range from $.8974$ to $.9009$: near-nominal components coexist with severe joint failure. Bonferroni instead gives $.9215$ [$ .9188,.9242$] and $.8981$ [$ .8950,.9013$]. The control fails the desired event, not its coordinatewise interpretation. Operative joint scores are near nominal in the core cells, but not every pointwise interval contains $.9$: interior $K=10$ L2 is $.8967$ [$ .8936,.8999$]. This small shortfall is below the $.01$ practical margin; many finite descriptive intervals do not refute the exchangeable rank theorem.

\paragraph{Q2: marginal coverage hides minority failure.}
Figure~\ref{fig:core_group_coverage} contrasts global, fitted-scale and fixed-group ILR. At $K=3$, global population coverage $.9008$ [$ .8980,.9036$] masks minority coverage $.1610$ [$ .1534,.1686$]. At $K=10$, population $.9003$ [$ .8975,.9030$] masks minority $.0511$ [$ .0371,.0652$]. Scaling raises minority coverage to $.8921$ [$ .8842,.8999$] and $.8885$ [$ .8808,.8961$]; group calibration gives $.9044$ [$ .8955,.9134$] and $.9096$ [$ .9008,.9184$]. Group-minus-global paired minority improvements are $.7435$ [$ .7317,.7553$] and $.8584$ [$ .8429,.8740$], exceeding the declared margin with mean half-widths below $.02$. Scaling helps here but its minority means remain below $.9$; only fixed-group calibration has the stated group theorem.

\paragraph{Q1--Q2: geometry ordering and adaptation costs.}
Figure~\ref{fig:geometry_reliability} places achieved joint coverage beside common area. Table~\ref{tab:principal_area_contrasts} uses paired statistical intervals enlarged by the separate numerical allowance. ILR-minus-L2 mean area is $-.04394$ [$-.04764,-.04024$] in the interior, but $+.04977$ [$ .04554,.05401$] near the boundary and $+.17705$ [$ .17023,.18388$] under the Dirichlet generator. These resolved reversals rule out a universal score ordering for this shared predictor. Every core-K3 TV/L2 comparison remains unresolved; interior/heterogeneous ILR versus Hellinger is also unresolved. A smaller naive-coordinate region with joint undercoverage is not better calibrated efficiency.
\input{figures/final_main}

In the heterogeneous $K=3$ cell, fitted-scale and group calibration increase population area over global ILR by $.01222$ [$ .00608,.01835$] and $.01384$ [$ .00767,.02002$], respectively (enlarged intervals). This measured cost accompanies substantially better minority allocation; it does not compare identical group guarantees. The scaled-versus-group area contrast is unresolved. The full audit covers 3,008 method/center regions and 3,616 paired numerical contrasts; the largest empirical allowance is $.004229$. Among 339 reported area contrasts, 76 remain unresolved and 258 meet the geometric practical-resolution flag. Those flags alone do not establish useful efficiency at the relevant coverage; the remaining sign-resolved small effects do not meet the practical margin. No K10 area is measured.

\paragraph{Q3: the specified shift and calibration scarcity.}
Changing only minority test prevalence from $.1$ to $.4$ reduces global ILR population coverage to $.6539$ [$ .6499,.6580$]. Fixed-group ILR gives $.9020$ [$ .8982,.9059$]; scaling gives $.8974$ [$ .8938,.9010$]. The complete within-group law stays fixed, so this agrees with the group-mixture derivation, not arbitrary-shift robustness (Figure~\ref{fig:prevalence_shift}). With total calibration size $99$, group ILR returns an infinite threshold/full domain for the minority in 31/100 trials; minority coverage is $.9493$ [$ .9364,.9622$]. This conservatism is retained, not discarded as a failure. At size $1999$, its minority mean is $.8943$ [$ .8882,.9003$], and population mean $.8978$ [$ .8959,.8998$]: more data does not force every finite-run mean above nominal.

For interior ILR, across-trial coverage SD decreases from $.03086$ at calibration size $99$ to $.01562$ at $499$ and $.01065$ at $1999$, with finite-test noise included. The size $99$ mean $.8939$ [$ .8877,.9000$] is another small pointwise shortfall. The $\alpha=.05$ interior arm gives $.9497$ [$ .9475,.9520$], targeting $.95$. Tables~\ref{tab:calibration_variability}--\ref{tab:small_group_cost} retain calibration-size/nominal effects and distinguish SD from mean uncertainty. Some scarce-group comparisons have minority mean half-width above $.02$; neither wide nor overlapping intervals establish equivalence.

\paragraph{Q3: zeros, smoothing and nearly full area.}
Unsmoothed ILR remains inapplicable to the entire actual-zero arm. The operative $\epsilon=.001$ score covers original compositions at $.8995$ [$ .8968,.9022$] for $K=3$ and $.8995$ [$ .8964,.9025$] for $K=10$. Yet $K=3$ zero-face coverage is $.5985$ [$ .5878,.6092$], positive-response coverage is $1$ in every trial, and mean full-dimensional area is $.997125$ [$ .996953,.997297$]. L2 zero-face coverage is $.6926$ [$ .6854,.6998$] at population coverage $.9017$ [$ .8988,.9046$]. Response-defined slices are descriptive, not test-observable groups with a conditional theorem.

Across $\epsilon=10^{-6},.001,.01$, smoothed-ILR face coverage stays near $.60$ (Table~\ref{tab:epsilon_sensitivity}). At $10^{-6}$, mean area is $1.000018$ [$ .999994,1.000041$], an estimate saturated near one. The audit retains 159 small out-of-range trial/group mean estimates (maximum $1.000374$) unclipped. These are numerical warnings, not true areas above one, proof of full-domain inclusion, or resolved cross-epsilon rankings. Faces have zero full-dimensional measure: near-full area cannot characterize their uncertainty. Full-domain bounds, unknown-status rates and infinite thresholds are reported separately.
